\protect \section *{Розділ~\ref *{sec:instability}. \nameref *{sec:instability}}
\begin{Solution}{8.{1}}
		$A=1/3$; $\gamma(1~\text{см})\approx46~\mathrm{с}^{-1}$, $\gamma(10~\text{см})\approx14{,}5~\mathrm{с}^{-1}$; швидше росте коротша хвиля (без поверхневого натягу). Умова $\omega^2=0$ дає $k_c$; при $k>k_c$ натяг переважає, $\omega^2>0$. У~плазмі роль натягу відіграє натяг силових ліній (збурення з~$k_\parallel\ne0$, шир) і~скінченний ларморівський радіус.
	
\end{Solution}
\begin{Solution}{8.{2}}
		$\omega_p\approx5{,}6\times10^{10}$~с$^{-1}$; $\gamma\approx0{,}069\,\omega_p\approx3{,}9\times10^{9}$~с$^{-1}$;
$\lambda=2\pi v_b/\omega_p\approx1{,}1$~мм; $t=10/\gamma\approx2{,}6$~нс, тобто $\approx23$ плазмові періоди.
	
\end{Solution}
\begin{Solution}{8.{3}}
		$\beta_\parallel\approx1{,}4$, $\beta_\perp\approx0{,}7$, різниця $\approx0{,}7<2$~--- плазма стійка.
	
\end{Solution}
\begin{Solution}{8.{4}}
		$t_{\text{зр}}=1/\gamma=5~\text{мкс}$ (амплітуда зростає в~$e$ разів); при $\gamma<0$ збурення загасає.
	
\end{Solution}
\begin{Solution}{8.{5}}
		межа нестійкості $\omega^2=0$ дає $k^2v_0^2=2\omega_p^2$; мінімум $\omega^2$ за~$k^2$ досягається при
$\sqrt{4k^2v_0^2+\omega_p^2}=2\omega_p$, де $\omega^2=-\omega_p^2/4$.
	
\end{Solution}
\begin{Solution}{8.{6}}
		$\delta W=mgl(1-\cos\theta)\approx mgl\,\theta^2/2>0$ біля $\theta=0$ (стійке); для $\theta=\pi+\vartheta$
$\delta W\approx-mgl\,\vartheta^2/2<0$ (нестійке). Оскільки $\omega^2_{\min}=\min(\delta W/K)$, одне зміщення з~$\delta W<0$ уже гарантує $\omega^2<0$.
	
\end{Solution}
\begin{Solution}{8.{7}}
		$|\mathcal E|\approx2{,}1\times10^{-3}$~ст.В $\approx0{,}6$~В; $m\Delta\theta=90^\circ\Rightarrow m=3$.
	
\end{Solution}
